\documentclass[12pt, a4paper]{article}

% ── Packages ──────────────────────────────────────────────────────────────────
\usepackage[margin=1in]{geometry}
\usepackage{times}             % Times New Roman font
\usepackage{setspace}          % Line spacing
\usepackage{titlesec}          % Title formatting
\usepackage{hyperref}          % Clickable links
\usepackage{xcolor}            % Colours
\usepackage{parskip}           % Paragraph spacing
\usepackage{microtype}         % Better typography
\usepackage{enumitem}          % List customisation

% ── Hyperlink styling ─────────────────────────────────────────────────────────
\hypersetup{
    colorlinks = true,
    urlcolor   = blue,
    linkcolor  = black,
    citecolor  = black
}

% ── Title customisation ───────────────────────────────────────────────────────
\titleformat{\section}{\large\bfseries}{}{0em}{}[\titlerule]
\titlespacing{\section}{0pt}{12pt}{6pt}

% ═════════════════════════════════════════════════════════════════════════════
\begin{document}

% ── Header ───────────────────────────────────────────────────────────────────
\begin{center}
    {\small Dhanani School Undergraduate Research Symposium (DURS) 2026}\\[4pt]
    {\small \textit{Stream: Artificial Intelligence \& Machine Learning}}\\[4pt]
    {\small \textit{Sub-stream: AI for Transforming Agriculture}}\\[18pt]

    {\LARGE \bfseries PQNK Knowledge Intelligence System: A RAG-Powered Expert\\[4pt]
    Chatbot for Sustainable Natural Farming in Pakistan}\\[14pt]

    {\normalsize
        \textbf{[Author 1 Name]}, \textbf{[Author 2 Name]}, \textbf{[Author 3 Name]}, \textbf{[Author 4 Name]}\\[4pt]
        \textit{Supervisor: [Supervisor Name]}\\[4pt]
        \textit{Dhanani School of Science \& Engineering, Habib University}\\[4pt]
        \textit{Industry Partner: Pakistan Agriculture Research (PAR)}\\[4pt]
    }
\end{center}

\vspace{6pt}
\hrule
\vspace{14pt}

% ── Abstract ─────────────────────────────────────────────────────────────────
\section*{Abstract}

\begin{onehalfspace}

In Pakistan, expert agronomic knowledge critical to smallholder farmers remains
fragmented across informal, heterogeneous media channels—including video lectures,
WhatsApp broadcasts, social media posts, and printed documents—rendering it
discoverable and inaccessible at scale. This research addresses that gap through the
\textbf{PQNK Knowledge Intelligence System}, an AI-powered knowledge repository
and expert conversational agent developed in collaboration with
\textbf{Pakistan Agriculture Research (PAR)}.

The system is centred on \textbf{Paedar Qudratti Nizam-e-Kashtari (PQNK)}, a
sustainable natural farming methodology pioneered by Dr.~Asif Sharif, whose domain
expertise currently reaches farmers exclusively through direct, manual consultation.
Our goal is to automate and scale that consultation process without compromising
the integrity of the source knowledge.

The system delivers two core contributions. First, we construct a
\textbf{multimodal knowledge base} by systematically collecting and processing the
scattered PQNK corpus. Raw content—spanning multiple languages and formats—undergoes
language detection, machine translation, semantic chunking, and vector embedding,
and is indexed into a Qdrant vector store for efficient retrieval. A structured web
application then exposes this unified repository, enabling farmers and researchers
to perform natural-language search across the entire knowledge base from a
single interface.

Second, we deploy a \textbf{domain-restricted Retrieval-Augmented Generation (RAG)
chatbot} that grounds every response exclusively within Dr.~Sharif's curated,
approved corpus, eliminating the hallucination risks of general-purpose LLMs.
Responses are delivered in both \textbf{text and synthesised audio} in Urdu and
English, deliberately lowering literacy and accessibility barriers for rural users.

The full-stack system integrates a React frontend with a Flask REST API, PostgreSQL
for user management, and Qdrant as the vector database. Expected contributions
include a replicable pipeline for expert knowledge digitalisation and a scalable
model for AI-driven agricultural extension services in multilingual, low-resource
settings.

\end{onehalfspace}

\vspace{12pt}

% ── Keywords ─────────────────────────────────────────────────────────────────
\noindent\textbf{Keywords:} Retrieval-Augmented Generation (RAG), Expert Knowledge Systems, Agricultural AI, Sustainable Farming, Natural Language Processing, Multimodal Knowledge Base, Large Language Models, Urdu NLP, Pakistan Agriculture.

\vspace{16pt}

\end{document}
